\chapter{บทสรุป}
\label{chapter:conclusion}

\section{สรุปผลการดำเนินงาน}

โครงงานนี้ได้พัฒนา UniAssist AI ซึ่งเป็นเว็บแอปพลิเคชันผู้ช่วยด้านวิชาการสำหรับนักศึกษาและอาจารย์ที่ปรึกษา คณะเทคโนโลยีสารสนเทศ สถาบันเทคโนโลยีพระจอมเกล้าเจ้าคุณทหารลาดกระบัง โดยมีเป้าหมายเพื่อรวบรวมข้อมูลทางการศึกษาที่กระจัดกระจายให้อยู่ในศูนย์กลางเดียว และให้บริการตอบคำถามด้านวิชาการได้อย่างถูกต้องและรวดเร็ว

ผลการดำเนินงานสามารถสรุปตามวัตถุประสงค์ได้ดังนี้

\begin{enumerate}
    \item \textbf{การรวบรวมและสร้างฐานข้อมูล} ระบบสามารถสกัดข้อมูลหลักสูตรจากเอกสาร มคอ.2 ด้วยโมเดลภาษาขนาดใหญ่ และจัดเก็บลงในฐานข้อมูลเชิงสัมพันธ์ที่ครอบคลุมทั้ง 4 หลักสูตร โดยมีการประเมินคุณภาพการสกัดเทียบกับเอกสารจริง และเลือกใช้โมเดล Qwen ที่ให้ผลสมดุลที่สุด (Coverage 91.7\%, Accuracy 95.2\%)
    \item \textbf{การพัฒนา AI Chatbot} ระบบใช้เทคนิค Text-to-SQL ในการแปลงคำถามภาษาไทยเป็นคำสั่ง SQL และตอบคำถามโดยอ้างอิงจากข้อมูลจริงในฐานข้อมูล พร้อมกลไกตรวจสอบความปลอดภัยของคำสั่งและกลไกสำรอง ทำให้คำตอบสามารถตรวจสอบย้อนกลับได้และลดปัญหาการสร้างข้อมูลเท็จ
    \item \textbf{การคำนวณเกรดและประเมินความเสี่ยง} ระบบสามารถคำนวณ GPA และประเมินสถานะทางการเรียนโดยอัตโนมัติ โดยอ้างอิงค่าเกรดและเกณฑ์จากฐานข้อมูล ไม่ใช้ค่าคงที่
    \item \textbf{ระบบทุนการศึกษาและแดชบอร์ดอาจารย์} ระบบให้ข้อมูลทุนการศึกษาและช่วยวิเคราะห์คุณสมบัติเบื้องต้น พร้อมแดชบอร์ดที่ช่วยให้อาจารย์ที่ปรึกษาติดตามภาพรวมและความเสี่ยงของนักศึกษาในความดูแลได้
    \item \textbf{การยืนยันตัวตนและความปลอดภัย} ระบบมีการยืนยันตัวตนผ่าน Google OAuth และแบ่งสิทธิ์ตามบทบาท ควบคู่กับการจำกัดการเข้าถึงฐานข้อมูลแบบอ่านอย่างเดียว
\end{enumerate}

\section{อภิปรายผล}

ผลการดำเนินงานแสดงให้เห็นว่าการเลือกวิธีการค้นคืนข้อมูลควรพิจารณาจากลักษณะของข้อมูลเป็นสำคัญ เนื่องจากข้อมูลหลักสูตรมีโครงสร้างที่ชัดเจน การใช้เทคนิค Text-to-SQL บนฐานข้อมูลเชิงสัมพันธ์จึงให้คำตอบที่แม่นยำ ตรวจสอบย้อนกลับได้ และลดปัญหาการสร้างข้อมูลเท็จได้ดีกว่าการให้โมเดลภาษาตอบจากความจำโดยตรง อีกทั้งการออกแบบให้โมเดลภาษาเป็นเพียงผู้สร้างคำสั่งค้นข้อมูล ไม่ใช่ผู้เก็บข้อเท็จจริง ทำให้ระบบมีความน่าเชื่อถือมากขึ้น

นอกจากนี้ ผลการประเมินการสกัดข้อมูลยังชี้ให้เห็นว่าไม่มีโมเดลใดที่ดีที่สุดในทุกมิติ การพิจารณาทั้งความครบถ้วนและความถูกต้องควบคู่กันจึงมีความสำคัญในการเลือกโมเดลให้เหมาะกับงาน

\section{ข้อจำกัดของระบบ}

\begin{enumerate}
    \item \textbf{ความไม่คงเส้นคงวาของบริการโมเดลภาษา} การเรียกใช้โมเดลผ่านผู้ให้บริการภายนอกอาจให้ผลลัพธ์ที่ไม่นิ่งเป็นบางครั้ง (เช่น ตอบ \texttt{NO\_ANSWER} แม้ตอบได้จริง) ซึ่งบรรเทาได้ด้วยกลไกลองใหม่และการเลือกโมเดลที่นิ่งกว่า แต่ยังไม่สามารถขจัดได้ทั้งหมด
    \item \textbf{ความครบถ้วนของการสกัดข้อมูล} รายวิชาบางประเภท เช่น วิชาเลือกเสรีหรือวิชาที่ใช้รหัส placeholder ในเอกสาร มคอ.2 ยังถูกสกัดได้ไม่ครบถ้วนสมบูรณ์
    \item \textbf{ขอบเขตของข้อมูล} ระบบครอบคลุมเฉพาะข้อมูลหลักสูตรของคณะเทคโนโลยีสารสนเทศ และข้อมูลนักศึกษาตัวอย่างที่ใช้ทดสอบยังมีจำนวนจำกัด
    \item \textbf{การค้นคืนข้อมูลเชิงบรรยาย} ระเบียบหรือประกาศที่อยู่ในรูปแบบข้อความยาวและไม่มีโครงสร้างชัดเจน ยังไม่ได้รับการค้นคืนด้วยเทคนิคเชิงความหมาย (Semantic Search) อย่างเต็มรูปแบบ
\end{enumerate}

\section{ข้อเสนอแนะและแนวทางการพัฒนาต่อ}

\begin{enumerate}
    \item เพิ่มการค้นคืนข้อมูลแบบ Retrieval-Augmented Generation ด้วยฐานข้อมูลเชิงเวกเตอร์ สำหรับข้อมูลเชิงบรรยายที่ไม่มีโครงสร้างชัดเจน เพื่อเสริมการทำงานของแกน Text-to-SQL ที่เหมาะกับข้อมูลเชิงโครงสร้าง
    \item เชื่อมต่อกับระบบทะเบียนของสถาบันโดยตรง เพื่อให้ข้อมูลผลการเรียนและตารางเรียนของนักศึกษาเป็นปัจจุบันแบบอัตโนมัติ
    \item ขยายฟังก์ชันการแนะนำรายวิชาและการตรวจสอบตารางเรียน (Schedule Conflict) ให้สมบูรณ์ยิ่งขึ้น
    \item ขยายชุดข้อมูลนักศึกษาและจัดทำการประเมินประสิทธิภาพการตอบคำถามของระบบกับผู้ใช้จริงในวงกว้าง
    \item พิจารณาปรับจูนโมเดลภาษา (Fine-tuning) ด้วยชุดคำถาม-คำสั่ง SQL เฉพาะโดเมนของสถาบัน เพื่อเพิ่มความแม่นยำและลดการพึ่งพาบริการภายนอก
\end{enumerate}

\section{สรุป}

โดยรวม UniAssist AI สามารถบรรลุวัตถุประสงค์ที่ตั้งไว้ในการเป็นเครื่องมือช่วยให้นักศึกษาเข้าถึงข้อมูลและวางแผนการเรียนได้อย่างมีประสิทธิภาพ ลดความเสี่ยงในการพ้นสภาพนักศึกษา และช่วยให้อาจารย์ที่ปรึกษาติดตามดูแลนักศึกษาได้อย่างทั่วถึง ซึ่งสอดคล้องกับเป้าหมายการพัฒนาที่ยั่งยืน SDG 4 ในการส่งเสริมโอกาสการเรียนรู้และคุณภาพการศึกษาอย่างเท่าเทียม
